% Виртуальная антенная решётка в 3D: математическое описание по статье
% Matveev, Konovalov — Distributed Continuous Time Algorithms... (preprint, 2023).
% Сборка: pdflatex grid_antenna.tex
\documentclass[12pt,a4paper]{article}

\usepackage[utf8]{inputenc}
\usepackage[T2A]{fontenc}
\usepackage[russian]{babel}
\usepackage{amsmath,amssymb,amsthm}
\usepackage{geometry}
\geometry{margin=2.5cm}
\usepackage{hyperref}
\hypersetup{colorlinks=true,linkcolor=blue,urlcolor=blue}

\title{Построение виртуальной антенной решётки в~3D:\\
координатная композиция протоколов (по разд.~7.4 статьи Matveev--Konovalov)}
\author{Текст восстановлен по математическому изложению в \texttt{BooksDocuments/nnrautom.pdf}.\\
Реализация по пути \texttt{Anatoliy/swarm\_mavic/grid\_antenna} в данной рабочей копии \emph{отсутствует}.}
\date{}

\begin{document}
\maketitle

\begin{abstract}
Заметка даёт сжатое изложение протокола построения вертикальной «виртуальной антенны'' из однотипных мобильных агентов над заданной точкой маяка на плоскости: горизонтальные координаты сходятся к проекции маяка, по вертикали агенты выстраиваются столбиком с шагом~$w$ над высотой~$h^\dagger$. Описание опирается на непрерывные одномерные протоколы равномерного развертывания и консенсуса из той же работы и на \emph{протокол~5} (разд.~7.4), который собирает трёхмерную скорость из покомпонентных скалярных законов с разнесением бюджета скорости и дальности видимости по осям. Без ссылок на код: только протоколы и условия сходимости в формулировках статьи.
\end{abstract}

\paragraph{Источник и ограничения.}
Ниже используются обозначения и шаги из препринта A.\,S.~Matveev, P.\,A.~Konovalov (\texttt{nnrautom.pdf}): разд.~2 (постановка, информанты, граф~$\Gamma$), разд.~3--5 (протоколы~1, \textbf{1max},~2), разд.~7.3--7.4 и \textbf{протокол~5}. Утверждения о предельной конфигурации формулируются в объёме \emph{предложения~2} и доказательственной схемы статьи (связность~$\Gamma_{\mathrm{ess}}$, выполнение дополнительных геометрических условий на параметры и начальное расположение). Имплементация под именем \texttt{grid\_antenna} в репозитории не проверялась.

%---------------------------------------------------------------------
\section{Постановка задачи}

Требуется расположить $N\ge 2$ одинаковых агентов в пространстве над горизонтальной плоскостью~$H$. У каждого агента есть приёмопередатчик; цель --- выстроить их \emph{вертикальной цепочкой} над выбранной точкой-маяком $p\in H$: расстояние между соседними агентами по вертикали должно стремиться к заданному шагу~$w>0$, а высота \emph{нижнего} агента --- к заданной отметке $h^\dagger>0$; горизонтальные координаты (в выбранной системе отсчёта с началом в~$p$) должны сходиться к нулю, т.\,е. проекции на~$H$ --- к точке~$p$.

Агенты анонимны, число~$N$ им неизвестно; доступны лишь \emph{относительные} координаты «видимых'' объектов (других агентов, маяка, плоскости) в пределах ограничений дальности и условий «прямой видимости'' в смысле статьи (см. предположение~1 и определение информантов в исходнике). Управление в одномерных протоколах --- скаляр $u_i$ с $|u_i|\le u$; в трёхмерном сценарии --- вектор скорости~$\mathbf{v}_i$ с евклидовым ограничением~$\|\mathbf{v}_i\|\le v$.

%---------------------------------------------------------------------
\section{Система координат и обозначения}

Вводится абсолютная декартова система~(AF) с началом в точке~$p$, осью~$z$, направленной перпендикулярно~$H$ «вверх'', и осями $x$, $y$ в плоскости~$H$. Положение агента~$i$ обозначим $\mathbf{r}_i(t)=(x_i(t),y_i(t),z_i(t))^\top$; $h_i$ --- его высота над~$H$. В этой рамке координата $z_i$ и высота $h_i$ совпадают по смыслу, поэтому далее удобно пользоваться обозначением $z_i$.

Пусть $R_{\mathrm{sen}}>0$ --- радиус сенсорного обнаружения в смысле статьи: при необходимости дальности видимости по осям $R_{\mathrm{vis}}^x,R_{\mathrm{vis}}^y,R_{\mathrm{vis}}^z$ подбирают так, что
\begin{equation}
  (R_{\mathrm{vis}}^x)^2+(R_{\mathrm{vis}}^y)^2+(R_{\mathrm{vis}}^z)^2 \le R_{\mathrm{sen}}^2,
\end{equation}
а для устойчивого наблюдения вблизи целевой конфигурации в разд.~7.4 предполагается $R_{\mathrm{sen}}>w$ и $R_{\mathrm{sen}}>h^\dagger$, и выбирают $R_{\mathrm{vis}}^z>w,h^\dagger$.

Для скорости задают положительные «бюджеты'' по осям $v_x,v_y,v_z$ так, что
\begin{equation}
  v_x^2+v_y^2+v_z^2 \le v^2,
\end{equation}
где $v$ --- максимальная допустимая евклидова скорость агента (см. ограничение на $\mathbf{v}_i$ ниже). На каждой оси выбирают отображение $\Xi_b:[0,\infty)\to[0,v_b]$ ($b\in\{x,y,z\}$), удовлетворяющее условиям типа~(7) статьи на соответствующем интервале видимости (см. следующий раздел).

%---------------------------------------------------------------------
\section{Базовые одномерные протоколы}

\subsection{Информанты и «существенные'' соседи}

У агента~$i$ в момент~$t$ множество информантов $I_i(t)$ задаётся правилами статьи: доступны относительные величины $x_j(t)-x_i(t)$ для одномерной координаты~$x$ (аналогично для других осей). Множество \emph{существенных} верхних/нижних информантов для координаты~$x$:
\begin{equation}
  I_i^{\pm}(t):=\Bigl\{j\in I_i(t): R_{\mathrm{vis}}>\pm\bigl(x_j(t)-x_i(t)\bigr)>0\Bigr\}.
\end{equation}

\subsection{Протокол~1: равномерное развертывание с шагом~$w$}

Положим $\Delta_{i,j}(t):=x_j(t)-x_i(t)$. Шаги:
\begin{enumerate}
  \item вычислить $I_i^{\pm}(t)$ по формуле выше;
  \item задать
  \begin{equation}
    d_i^{\pm}(t):=
    \begin{cases}
      \displaystyle\min_{j\in I_i^{\pm}(t)} |\Delta_{i,j}(t)|, & I_i^{\pm}(t)\neq\varnothing,\\[0.7em]
      w, & I_i^{\pm}(t)=\varnothing;
    \end{cases}
  \end{equation}
  \item управление
  \begin{equation}
    u_i(t):=\Xi\bigl(d_i^{+}(t)\bigr)-\Xi\bigl(d_i^{-}(t)\bigr),
  \end{equation}
  где $\Xi:[0,\infty)\to[0,u]$ гладко на $[0,R_{\mathrm{vis}})$, $\Xi'(d)>0$ там, где требуется, и $\Xi(0)=0$ (точные условия --- формула~(7) в статье).
\end{enumerate}
Интерпретация: $d_i^{+}$ и $d_i^{-}$ --- расстояния до ближайших соседей «выше'' и «ниже'' по координате среди существенных информантов; если соседа нет, величина~$w$ задаёт «эталонный'' зазор. Разность $\Xi(d_i^{+})-\Xi(d_i^{-})$ задаёт дрейф к выравниванию зазоров и, в пределе, к решётке с шагом~$w$ (теорема~3.1 статьи при выполнении предпосылок).

\subsection{Протокол \textbf{1max}: консенсус и консенсус с маяком}

Модификация для задач согласования: в~(5) вместо $\min$ берётся $\max$ по множеству $I_i^{\pm}$, \emph{параметр $w$ полагается равным нулю}, а в определении существенных соседей и в условиях на~$\Xi$ формально допускается $R_{\mathrm{vis}}:=\infty$, что расширяет множество информантов (формулировка в разд.~4 статьи). Поэтому в режиме \textbf{1max}, если $I_i^{\pm}(t)=\varnothing$, соответствующая величина $d_i^{\pm}(t)$ равна нулю, а не шагу развертывания. Управление по-прежнему $u_i=\Xi(d_i^{+})-\Xi(d_i^{-})$. Теорема~4.1: при симметрии связей и связности~$\Gamma_{\mathrm{ess}}$ получается консенсус; в варианте с якорем $x^\dagger$ --- сходимость к~$x^\dagger$.

В задаче антенны по осям $x$ и $y$ маяк $p$ становится «фиктивным агентом'' с координатами $0$ в AF; протокол~\textbf{1max} применяют к относительным координатам всех видимых объектов, включая~$p$ при доступности, что ведёт к сходимости $x_i,y_i\to 0$ (т.\,е. к проекции на~$p$) --- как в доказательстве предложения~2.

\subsection{Протокол~2: развертывание с якорем $x^\dagger$ и шагом~$w$}

Каждый агент выполняет протокол~1, считая $x^\dagger$ состоянием «товарища'', если якорь доступен как информант; \emph{если множество нижних существенных информантов пусто}, во второй строке формулы~(5) вместо~$w$ подставляют $R_{\mathrm{vis}}$ (текст статьи после протокола~2). Теорема~5.1: при выполнении предпосылок достигается равномерное развертывание с шагом~$w$ относительно якоря в смысле п.~4 определения~2.1 статьи (нижний агент у якоря, остальные через $w$, $2w$, \ldots\ в перечислении).

%---------------------------------------------------------------------
\section{Трёхмерный протокол~(протокол~5, разд.~7.4)}

\subsection{Кинематика и ограничение на скорость}

Агенты описываются интегратором
\begin{equation}
  \dot{\mathbf{r}}_i=\mathbf{v}_i,\qquad \|\mathbf{v}_i\|\le v,
\end{equation}
что в статье обозначено как ограничение типа~(8) для векторного управления (в 2D перед разд.~7.3 явно указано $\|\mathbf{v}_i\|\le v$; в разд.~7.4 компоненты собираются в трёхмерный вектор).

\subsection{Покомпонентное построение $\mathbf{v}_i=(v_i^x,v_i^y,v_i^z)^\top$}

\textbf{Шаг 1 ($b\in\{x,y\}$).} Вычислить скаляр $v_i^b(t)$ как выход протокола~\textbf{1max} с картой $\Xi:=\Xi_b$, радиусом существенных связей $R_{\mathrm{vis}}^b$ и относительными координатами $b_j(t)-b_i(t)$ по всем видимым объектам~$j$, \emph{включая маяк}~$p$ (координаты $0$ в AF), если относительное положение маяка доступно.

\textbf{Шаг 2 ($z$).} Вычислить $v_i^z(t)$ по протоколу~2 с $\Xi:=\Xi_z$, $R_{\mathrm{vis}}:=R_{\mathrm{vis}}^z$, якорем $x^\dagger:=h^\dagger$ и шагом~$w$, применённым к относительным аппликатам $z_j(t)-z_i(t)$ видимых объектов.

\textbf{Шаг 3.} Положить
\begin{equation}
  \mathbf{v}_i(t)=\bigl(v_i^x(t),\,v_i^y(t),\,v_i^z(t)\bigr)^\top.
\end{equation}

Явно, через те же промежутки $d_{i,b}^{\pm}(t)$, что в протоколах~1/\textbf{1max}/2 (но для каждой оси свои $I_i^{\pm}$, свои $R_{\mathrm{vis}}^b$ и для~\textbf{1max} по $x,y$ --- вариант с $\max$ в определении $d_{i,b}^{\pm}$):
\begin{align}
  v_i^x(t) &= \Xi_x\!\bigl(d_{i,x}^{+}(t)\bigr)-\Xi_x\!\bigl(d_{i,x}^{-}(t)\bigr),\\
  v_i^y(t) &= \Xi_y\!\bigl(d_{i,y}^{+}(t)\bigr)-\Xi_y\!\bigl(d_{i,y}^{-}(t)\bigr),\\
  v_i^z(t) &= \Xi_z\!\bigl(d_{i,z}^{+}(t)\bigr)-\Xi_z\!\bigl(d_{i,z}^{-}(t)\bigr),
\end{align}
где для $b=x,y$ величины $d_{i,b}^{\pm}$ вычисляются по правилам~\textbf{1max} на координате~$b$ (с учётом маяка как объекта с $b=0$), а для $b=z$ --- по протоколу~2 с якорем $h^\dagger$ и шагом~$w$ (с подстановкой $R_{\mathrm{vis}}^z$ вместо~$w$ при $I_i^{-}=\varnothing$ на оси~$z$).

В статье отмечается, что при выборе $\Xi_b$ с значениями в $[0,v_b]$ и выполнении $v_x^2+v_y^2+v_z^2\le v^2$ ограничение $\|\mathbf{v}_i\|\le v$ сохраняется (аналогично оценке~(9) для 2D-композиции).

\subsection{Роль $R_{\mathrm{vis}}^x,R_{\mathrm{vis}}^y,R_{\mathrm{vis}}^z$ и бюджетов $v_x,v_y,v_z$}

Величины $R_{\mathrm{vis}}^b$ задают, на каких масштабах по каждой оси выбираются «существенные'' соседи для протоколов; их квадраты суммируются в ограничение на $R_{\mathrm{sen}}$ (см. разд.~2 настоящей заметки), чтобы локальные одномерные измерения укладывались в глобальный радиус восприятия. Бюджеты $v_b$ распределяют максимально допустимые модули покомпонентных вкладов так, чтобы евклидова норма итоговой скорости не превосходила~$v$.

\subsection{Целевая конфигурация (по предложению~2)}

В статье для сценария с протоколом~5 при указанных начальных условиях (агенты в разных высотах $h_i(0)>h^\dagger$ над прямоугольником $|x|<R_{\mathrm{vis}}^x$, $|y|<R_{\mathrm{vis}}^y$ при $z=0$, плюс прочие предпосылки теорем о протоколах) формулируется: роботы не сталкиваются друг с другом и с~$H$, проекции на~$H$ сходятся к~$p$, и существует перечисление агентов, для которого $z_i(t)\to h^\dagger+i w$ при $t\to\infty$. Горизонтально это соответствует $x_i,y_i\to 0$ в AF. Важная особенность формулировки предложения~2: исходный граф информационных связей $\Gamma_{\mathrm{all}}(0)$ может быть вообще \emph{пустым} (без рёбер), например из-за больших начальных разрывов по высоте; полная картина доказательства опирается на последовательное применение теорем из разд.~4--5 и оценки видимости из доказательства.

%---------------------------------------------------------------------
\section{Поэтапная идея алгоритма}

\begin{enumerate}
  \item Зафиксировать маяк $p$ и ввести систему AF с началом в~$p$; задать $w$, $h^\dagger$, $v$, $R_{\mathrm{sen}}$, а также допустимые $R_{\mathrm{vis}}^x,R_{\mathrm{vis}}^y,R_{\mathrm{vis}}^z$ и $v_x,v_y,v_z$ с условиями сумм квадратов.
  \item Для каждого агента на шаге~$t$ измерить относительные координаты видимых объектов (агенты, маяк, плоскость~$H$ в той мере, в какой это заложено в модель сенсоров статьи).
  \item По координатам $x$ и $y$ применить протокол~\textbf{1max} с функциями $\Xi_x,\Xi_y$, получив $v_i^x,v_i^y$ как стремление к консенсусу в нуле (проекция на~$p$).
  \item По координате $z$ применить протокол~2 с якорем $h^\dagger$ и шагом~$w$, получив $v_i^z$ как вертикальное равномерное развертывание «над'' якорной высотой.
  \item Собрать $\mathbf{v}_i=(v_i^x,v_i^y,v_i^z)^\top$, соблюдая $\|\mathbf{v}_i\|\le v$.
  \item В пределе при выполнении предпосылок статьи: горизонтальное сжатие к~$p$ и вертикальная решётка $z_i\to h^\dagger+iw$ при подходящей нумерации.
\end{enumerate}

%---------------------------------------------------------------------
\section*{Простыми словами}

Каждый робот смотрит, кто у него «сверху'' и «снизу'' по высоте и кто слева-справа в горизонтальной плоскости относительно точки маяка, но только среди тех, кого он реально «видит'' на ограниченной дистанции. По горизонтали он слегка подстраивается так, чтобы согласовать свою позицию с остальными и с маяком --- все съезжаются над одной точкой. По вертикали тот же тип правила выравнивает интервалы между соседями и привязывает нижний этаж к заданной высоте, так что колонна превращается в ровную лестницу с одинаковой ступенькой~$w$. Ограничение на скорость распределяют по осям, чтобы итоговый вектор движения не был слишком быстрым.

\end{document}
